\documentclass[10pt]{article}
\usepackage[T1]{fontenc}
\usepackage[utf8]{inputenc}
\usepackage[french]{babel}
\usepackage{amsmath,amssymb,mathrsfs,mathpazo}
\usepackage{geometry}
\geometry{a4paper,margin=2cm}
\usepackage{graphicx,array,tabularx,tabularray,colortbl,multirow,multicol,booktabs,enumitem,pifont,ulem,fancybox,diagbox,cancel}
\usepackage{pstricks,pstricks-add,pst-plot,pst-tree,pst-node,pst-text,pst-func,pst-math,pst-bspline,pst-3dplot,pst-eucl,pst-solides3d}
\let\pscalculate\undefined
\usepackage{pst-bezier}
\usepackage{tikz,pgfplots,tkz-tab}
\usetikzlibrary{arrows,arrows.meta,patterns,calc,positioning,trees,decorations.markings,shapes.geometric,shapes.misc,shapes.arrows,intersections,decorations.pathmorphing,decorations.pathreplacing,backgrounds,shadows,babel}
\pgfplotsset{compat=1.18}
\usepgfplotslibrary{fillbetween}
\usepackage[np]{numprint}
\usepackage{listings,makecell,pas-tableur,twoopt}
\let\barre\undefined
\usepackage[pstricks]{tablvar}
\definecolor{darkgreen}{RGB}{0,100,0}
\providecommand{\red}{\color{red}}
\providecommand{\blue}{\color{blue}}
\providecommand{\green}{\color{green}}
\providecommand{\black}{\color{black}}
\providecommand{\white}{\color{white}}
\providecommand{\gray}{\color{gray}}
\providecommand{\cyan}{\color{cyan}}
\providecommand{\magenta}{\color{magenta}}
\providecommand{\yellow}{\color{yellow}}
\providecommand{\violet}{\color{violet}}
\providecommand{\orange}{\color{orange}}
\providecommand{\euro}{\texteuro}
\providecommand{\EUR}{\texteuro}
\providecommand{\ieme}{\textsuperscript{e}}
\providecommand{\degres}{\ensuremath{^\circ}}
\providecommand{\degree}{\ensuremath{^\circ}}
\providecommand{\vv}[1]{\overrightarrow{#1}}
\providecommand{\Coord}[1]{#1}
\providecommand{\decosix}{}
\providecommand{\starredbullet}{\ensuremath{\bullet}}
\pagestyle{empty}
\setlength{\parindent}{0pt}
\usepackage{adjustbox}\DeclareUnicodeCharacter{2236}{:}\DeclareUnicodeCharacter{2190}{\ensuremath{\leftarrow}}
\providecommand{\R}{\mathbb{R}}
\providecommand{\N}{\mathbb{N}}
\providecommand{\Z}{\mathbb{Z}}
\providecommand{\Q}{\mathbb{Q}}
\providecommand{\D}{\mathbb{D}}
\providecommand{\C}{\mathbb{C}}
\providecommand{\e}{{\mathrm{e}}}
\providecommand{\vect}[1]{\overrightarrow{#1}}
\providecommand{\vectt}[1]{\overrightarrow{\mathrm{#1}}}
\providecommand{\vvt}[1]{\overrightarrow{\mathrm{#1}}}
\providecommand{\np}[2][]{{#2~#1}}
\providecommand{\vv}[1]{{\overrightarrow{#1}}}
\renewcommand{\d}{{\mathrm{d}}}
\providecommand{\degres}{{^\circ}}
\providecommand{\degre}{{^\circ}}
\providecommand{\blue}{\color{blue}}
\providecommand{\red}{\color{red}}
\providecommand{\enskip}{\quad}
\providecommand{\llbracket}{\lbrack\!\lbrack}
\providecommand{\rrbracket}{\rbrack\!\rbrack}
\providecommand{\euro}{{\text{€}}}
\providecommand{\no}{{\text{n°}}}
\providecommand{\interval}[2]{{[#1;#2]}}
\providecommand{\intervFO}[2]{{[#1;#2[}}
\providecommand{\dotfill}{\ldots}

\providecommand{\ds}{\displaystyle}
\providecommand{\textding}[1]{#1}
\providecommand{\Oij}{\(\mathrm{O};\vec{\imath},\vec{\jmath}\)}
\providecommand{\Oijk}{\(\mathrm{O};\vec{\imath},\vec{\jmath},\vec{k}\)}
\providecommand{\Ouv}{\(\mathrm{O};\vec{u},\vec{v}\)}
\providecommand{\intervFF}[2]{\ensuremath{[#1;#2]}}
\renewcommand{\euro}{\ifmmode\text{\texteuro}\else\texteuro\fi}
\DeclareUnicodeCharacter{00B0}{\ifmmode^{\circ}\else\textdegree\fi}
\DeclareUnicodeCharacter{20AC}{\ifmmode\text{\texteuro}\else\texteuro\fi}
% Énoncés : Éducation nationale. Transcription et corrigés : APMEP, Denis Vergès et contributeurs. Découpage et classement : CoopMaths, Sébastien Lozano, Éric Elter et contributeurs. Adaptation de la présentation : Maths Explorer.
% CC BY-SA : https://coopmaths.fr/www/about/
% https://coopmaths.fr/alea/static/bac/2021/tex/bac_2021_06_etrangers_5_cor.tex
\begin{document}
\noindent\textit{APMEP / CoopMaths. CC BY-SA. Présentation adaptée par Maths Explorer.}\par\medskip
\bigskip 

\textbf{Partie A}

\medskip

Soit $g$ la fonction définie sur $\R$ par :
$g(x) = 2\e^{\frac{-1}{3}x} + \dfrac{2}{3}x - 2.$

\smallskip

\begin{enumerate}
\item On admet que la fonction $g$ est dérivable sur $\R$ et on note $g'$ sa fonction dérivée.


$g’(x) = 2\times \left ( -\dfrac{1}{3}\right ) \e^{\frac{-1}{3}x} + \dfrac{2}{3} = \dfrac{-2}{3}\text{e}^{- \frac{1}{3}x} + \dfrac{2}{3}$

\item 
$g'(x)=\dfrac{2}{3}\left ( 1- \e^{\frac{-1}{3}x}\right ) $

$g'(x)>0 
\iff 1- \e^{\frac{-1}{3}x} >0
\iff 1> \e^{\frac{-1}{3}x}
\iff \ln(1) > \dfrac{-1}{3}x
\iff 0 > \dfrac{-1}{3}x
\iff x > 0$

$g(0)=2\e^{0}- 0 - 2 = 0$ 

On en déduit les variations de la fonction $g$:

\begin{center}
















\psset{unit=1cm}
\begin{pspicture}(7,3)
\psframe(7,3)\psline(0,2)(7,2)\psline(0,2.5)(7,2.5)\psline(1,0)(1,3)
\psline{->}(1.5,1.5)(3.5,0.5)\psline{->}(4.5,0.5)(6.5,1.5)
\uput[u](0.5,2.4){$x$} \uput[u](1.5,2.4){$-\infty$} \uput[u](4,2.4){$0$} \uput[u](6.5,2.4){$+ \infty$} 
\uput[u](0.5,1.9){$g'(x)$} \uput[u](2.5,1.9){$-$} \uput[u](4,1.9){$0$} \uput[u](5.5,1.9){$+$} \uput[u](4,0){$0$}
\rput(0.5,1){$g(x)$}
\end{pspicture}
\end{center}	

\item 
D'après le tableau de variations de $g$, on a : $g(x)\geqslant 0$ pour tout réel $x$.
\end{enumerate}

Un tracé de la courbe représentative sur la calculatrice le confirme.
\begin{center}
\psset{unit=1cm}
\begin{pspicture}(-5,-1)(5,4)
\psaxes[linewidth=1.25pt,labelFontSize=\scriptstyle]{->}(0,0)(-5,0)(5,4)
\psplot[plotpoints=2000,linewidth=1.25pt,linecolor=red]{-4.5}{5}{2 x mul 3 div 2 sub 2.71828 x 3 div neg exp 2 mul add}
\end{pspicture}
\end{center}
\bigskip

\textbf{Partie B}

\medskip

\begin{enumerate}
\item On considère l'équation différentielle
$\qquad (E) :\qquad  3y' + y = 0.$



$(E) \iff y'+\dfrac{1}{3}y=0$ qui a pour solutions d'après le cours, les fonctions $h$ définies par $h(x)=C\times \e^{\frac{-1}{3}x}$, où $C\in\R$.

\item La solution particulière dont la courbe représentative, dans un repère du plan, passe par le point M(0~;~2) vérifie $h(0)=2$, c'est-à-dire $C\e^{0}=2$, donc $C=2$.

La solution particulière cherchée est la fonction $h$ définie par $h(x)=2\e^{\frac{-1}{3}x}$.

\item Soit $f$ la fonction définie sur $\R$ par :
$f(x) = 2\text{e}^{- \frac{1}{3}x}$
 et $\mathcal{C}_f$ sa courbe représentative.
 
	\begin{enumerate}
		\item La tangente $\left(\Delta_0\right)$ à la courbe $\mathcal{C}_f$ en M(0~;~2) a pour équation:
		
$y=f'(0)\left (x-0\right ) + f(0)$		

$f(x)=2\e^{\frac{-1}{3}x}$ donc $f(0)=2$

$f'(x)= 2\times \left ( -\dfrac{1}{3} \right ) \e^{\frac{-1}{3}x} = -\dfrac{2}{3}\e^{\frac{-1}{3}x}$ donc $f'(0)=-\dfrac{2}{3}$
		
$\left(\Delta_0\right)$ a donc pour équation $y= - \dfrac{2}{3}x + 2$.

		\item Pour étudier, sur $\R$, la position de cette courbe  $\mathcal{C}_f$ par rapport à la tangente $\left(\Delta_0\right)$, on étudie le signe de $f(x)-\left ( -\dfrac{2}{3}x +2\right )$.
		
$f(x)-\left ( -\dfrac{2}{3}x +2\right ) = 2\e^{\frac{-1}{3}x} +\dfrac{2}{3}x - 2 = g(x)$; on sait d'après la partie A, que $g(x)\leqslant 0$ sur $\R$, donc la courbe $\mathcal{C}_f$ est toujours en dessous de la tangente $\Delta_0$.		
		
	\end{enumerate}
\end{enumerate}

\bigskip

\textbf{Partie C}

\medskip

\begin{enumerate}

\item Soit A le point de la courbe $\mathcal{C}_f$ d'abscisse $a$,\, $a$ réel quelconque.

La tangente $\left(\Delta_a\right)$ à la courbe $\mathcal{C}_f$ au point A a pour équation
$y=f'(a) \left (x-a\right )+f(a)$ c'est-à-dire 
$y=-\dfrac{2}{3}\e^{\frac{-1}{3}x}\left (x-a\right )+2\e^{\frac{-1}{3}a}$.

Elle coupe l'axe des abscisses en un point P dont l'abscisse est solution de l'équation \\
$-\dfrac{2}{3}\e^{\frac{-1}{3}a}\left (x-a\right )+2\e^{\frac{-1}{3}a} = 0$.
On résout cette équation:

$\begin{gathered}-\dfrac{2}{3}\e^{\frac{-1}{3}a}\left (x-a\right )+2\e^{\frac{-1}{3}a} = 0
\iff  2\e^{\frac{-1}{3}a} = \dfrac{2}{3}\e^{\frac{-1}{3}a}\left (x-a\right )
\iff \dfrac{2\e^{\frac{-1}{3}a}}{\frac{2}{3}\e^{\frac{-1}{3}a}}=x-a\\
\phantom{-\dfrac{2}{3}\e^{\frac{-1}{3}a}\left (x-a\right )+2\e^{\frac{-1}{3}a} = 0}
\iff 3 = x-a \iff x=a+3\end{gathered}$

Donc $\left(\Delta_a\right)$ coupe l'axe des abscisses en un point P de coordonnées $(a + 3~;~0)$.

\item La tangente $\left(\Delta_{-2}\right)$  à la courbe $\mathcal{C}_f$ au point B d'abscisse $- 2$ coupe l'axe des abscisses au point P de coordonnées $(-2+3~;~0)$ soit $(1~;~0)$.

La tangente $\left(\Delta_{-2}\right)$ est donc la droite (BP).
\end{enumerate}

\end{document}
